\clearpage
\SourceLeaf{247}%
\Needspace{7\baselineskip}
\begin{center}
{\Large[\textsc{Congruence and Symmetry.}]}
\end{center}

\par\bigskip

\Needspace{7\baselineskip}
\begin{center}
[\textsc{Gerling} \emph{to} \textsc{Gauss}.\quad \emph{Cassel, 25 March 1813.}]
\end{center}

\{\mbox{.\ .\ .\ .\ .\ .} In the December issue of last year's Monatliche Correspondenz, Professor \textsc{Mollweide} declares the usual proof for the surface area of spherical triangles inadmissible. It treats as equal two triangles whose 3 corresponding sides are equal but which cannot be congruent because their sides are differently situated, and this proposition does not occur in the treatises on spherical trigonometry. I have found a proof of this proposition that seems very simple to me. It is as follows. Two such triangles always arise when the radii of the sphere that form one of them are extended beyond the center. If planes are then drawn through the vertices of the two triangles, it is very easy to see that they are parallel. It then follows just as easily that they are equally far from the center. --- They thus trace out congruent small circles and congruent segments of the surface on the sphere. The triangles result when one subtracts from these congruent surface segments the 3 digons bounded by equal arcs of the small circles and by the equal sides of the triangles. These digons, however, are congruent, as is readily seen from the construction, so the triangles are equal as well.
\mbox{.\ .\ .\ .\ .\ .\ .}\}

\par\bigskip
